\section{长答案怎样评估：专家与普通用户的双重视角}

前一章说明了为什么需要开放问答，本章解决更难的问题：一段没有唯一标准答案的医学文本，究竟怎样判断好坏？课程把 evaluation 拆成多个轴，并让 clinicians 与 laypeople 分别评价自己最有能力判断的部分。

\subsection{评价框架不是一个总分}

Evaluation framework 页是后续所有结果图的阅读钥匙。专家评价关注与 scientific consensus 的一致性、是否展示正确或错误的 comprehension、retrieval 与 reasoning、是否有 missing content、潜在 harm 与 bias；普通用户评价则更关注答案是否回应问题意图、是否 helpful。每个轴都对应不同失败模式，因此不能用一个平均分掩盖危险短板。

\lecturefigure{slide-19-evaluation-framework.jpg}{长答案的人类评价框架}{官方录像恢复幻灯片 V019；Stanford CS25 Lecture 17}

\lecturefigure{slide-20-clinician-evaluation-rubric.jpg}{临床专家评价维度}{官方录像恢复幻灯片 V020；Stanford CS25 Lecture 17}

\lecturefigure{slide-21-layperson-evaluation-rubric.jpg}{普通用户评价维度}{官方录像恢复幻灯片 V021；Stanford CS25 Lecture 17}

\begin{longtable}{@{}L{0.22\textwidth}L{0.32\textwidth}L{0.34\textwidth}@{}}
\toprule
评价轴 & 主要问题 & 典型失败 \\
\midrule
scientific consensus & 回答是否符合当时可靠医学知识？ & 给出与共识冲突或证据不足的断言 \\
comprehension / retrieval / reasoning & 是否理解问题、取回相关知识并正确组合？ & 同一回答同时包含正确与错误推断 \\
incorrect or missing content & 是否加入无关错误，或漏掉关键内容？ & 详细但跑题；忽略必须就医的危险信号 \\
extent / likelihood of harm & 错误可能造成多严重、多可能的伤害？ & 低估病情、误导延迟治疗、错误自我处置 \\
bias & 是否出现不当人口统计或社会偏见？ & 无证据地把群体属性与诊断或价值判断绑定 \\
intent / helpfulness & 是否回应提问者真正关心的问题？ & 技术上正确但无法理解、无法行动 \\
\bottomrule
\end{longtable}

\begin{importantbox}{医学评价的向量视角}
把每个回答记为评价向量 $\mathbf{e}=(c,r,m,h,b,u)$，分别表示共识、推理、缺失、伤害、偏见与效用。系统改进不是让一个总分上升，而是寻找不能接受的坐标：例如 helpfulness 提升但 harm 也上升，就不能称为净改进。
\end{importantbox}

这种向量评价还改变了模型迭代方式。假设新 prompt 让 scientific consensus 上升，却让 missing content 与 verbosity 同时增加，团队就需要检查新增段落，而不是宣布版本整体更好；如果 layperson helpfulness 上升但 clinician harm 评价恶化，则面向患者的可读性优化越过了安全边界。实际审计可以把每次模型变更写成“哪个维度改善、哪个维度退化、退化是否可接受、由谁签字”的 change log。这样，human evaluation 不再只是论文末尾的展示，而成为发布门禁。课堂的 rubric 因此也提供了一套通用工程模板：先把失败模式拆开，再给每个维度找到合格评审者和升级路径。

\subsection{从通用语言模型到医学域对齐}

对齐页提出下一步：基础模型已经拥有广泛语言和知识能力，如何让它适应医学问题的格式与评价标准？这里的“domain alignment”不只是追加专业语料，而是让模型学会在医学指令下选择相关知识、组织长答案，并尽可能满足上述多维 rubric。

\lecturefigure{slide-22-aligning-medical-domain.jpg}{将大型语言模型对齐到医学领域}{官方录像恢复幻灯片 V022；Stanford CS25 Lecture 17}

\teachervoice{讲者在问答中明确承认 human evaluation set 仍然很小，约一百四十道长问题。昂贵的专家标注既是研究限制，也是值得建设的公共资源；因此，后面图表应被理解为有价值但初步的人类证据，而不是覆盖全部临床情境的统计保证。}

\begin{warningbox}{评审者身份决定可回答的问题}
患者可以判断是否易懂和有帮助，却不一定能判断细节是否符合临床共识；医生能判断医学内容，却未必代表所有患者的沟通需求。双重评价不是互相替代，而是把不同 epistemic role 放到正确位置。
\end{warningbox}

\subsection{本章小结}

长答案评价必须拆成多个风险与效用维度，并区分医生和普通用户的判断职责。这个框架为后续 Med-PaLM 结果提供了坐标系，也说明临床模型研究的瓶颈不仅是训练数据，还包括高质量、可复核的人类评价数据。

